\documentclass[11pt, a4paper]{scrartcl}
\usepackage{azzam}

\begin{document}

\section*{Mini Simulation 3 : OSK SMP 2026 6 Mei}
\begin{enumerate}
    \item Jika $x$ bilangan real, nilai terkecil dari
    $$\left (x^2-x+1\right )^2 -\left (x^2-x+1\right )+1$$
    yang mungkin adalah . . . .
    \begin{enumerate}[label={(\Alph*)}]
        \item $\dfrac{4}{5}$ \item $\dfrac{13}{16}$ \item $\dfrac{4}{7}$ \item $\dfrac{15}{22}$
    \end{enumerate}

    \item Seekor semut awalnya berada pada posisi $(0, 0)$ pada bidang kartesius. Setiap detiknya si semut berpindah tepat $1$ unit ke kanan, atau ke kiri, atau ke atas, atau ke bawah. Peluang bahwa setelah $4$ detik si semut berada di $(0, 0)$ adalah . . . .
    \begin{enumerate}[label={(\Alph*)}]
        \item $\dfrac{39}{256}$ \item $\dfrac{19}{128}$ \item $\dfrac{9}{64}$ \item $\dfrac{1}{4}$
    \end{enumerate}

    \item Pada segi-dua belas beraturan $ABCDEFGHIJKL$ berlaku panjang $AG=20$. Jarak titik $C$ ke garis $GL$ adalah . . . .
    \begin{enumerate}[label={(\Alph*)}]
        \item $10\sqrt{2}$ \item 10 \item $5\sqrt{6}$ \item 5
    \end{enumerate}

    \item Banyaknya bilangan asli kurang dari atau sama dengan $100$ yang jumlah semua faktor positifnya merupakan bilangan ganjil adalah . . . .
    \begin{enumerate}[label={(\Alph*)}]
        \item 16 \item 17 \item 18 \item 19
    \end{enumerate}

    \item Diketahui bahwa di dalam suatu kelas yang berisi $30$ murid, seorang guru matematika membagi kelas ini ke dalam 3 kelompok. Kelompok pertama berisi $a\ge 1$ murid yang semua nilai ulangannya $80$. Kelompok kedua berisi $b\ge 1$ murid yang semua nilai ulangannya $90$. Sedangkan kelompok ketiga berisi $c\ge 1$ murid yang semua nilai ulangannya $100$. Diketahui bahwa rata-rata nilai ulangan secara keseluruhan di kelas tersebut adalah $85$. Banyaknya tripel $(a,b,c)$ yang mungkin adalah . . . .
    \begin{enumerate}[label={(\Alph*)}]
        \item 6 \item 7 \item 8 \item 9
    \end{enumerate}

    \item Diketahui
    $$ A =\frac{1}{1\times 2} + \frac{1}{3\times 4}+\frac{1}{4\times 5}+\cdots + \frac{1}{149\times 150} $$
    dan 
    $$ B =\frac{1}{76\times 150} + \frac{1}{77\times 149} + \frac{1}{78\times 148}+\cdots + \frac{1}{150\times 76}. $$
    Nilai dari $\frac{B}{A}$ adalah . . . .
    \begin{enumerate}[label={(\Alph*)}]
        \item $\dfrac{1}{111}$ \item $\dfrac{1}{113}$ \item $\dfrac{1}{115}$ \item $\dfrac{1}{117}$
    \end{enumerate}

    \item Diberikan sebuah papan $8\times 8$ yang diisi sebuah bilangan di setiap kotaknya, dengan aturan bahwa kotak di baris ke-$i$ dan kolom ke-$j$ diisi $i\times j$. Kemudian, semua kotak yang memuat bilangan genap akan diwarnai hitam. Jumlah dari semua bilangan pada kotak diwarnai hitam adalah . . . .
    \begin{enumerate}[label={(\Alph*)}]
        \item 400 \item 720 \item 800 \item 1040
    \end{enumerate}

    \item Di dalam persegi $ABCD$ terdapat dua titik $E$ dan $F$ yang memenuhi $\angle DEF=\angle EFB=60^\circ$. Jika panjang $DE=6$, $EF=10$, dan $FB=8$, maka luas dari persegi $ABCD$ adalah . . . .
    \begin{enumerate}[label={(\Alph*)}]
        \item 78 \item 84 \item 88 \item 62
    \end{enumerate}

    \item Banyaknya pasangan bilangan asli $(a,b)$ yang memenuhi
    $$\frac{a}{b} + \frac{1-2b}{a} =4$$
    adalah . . . .
    \begin{enumerate}[label={(\Alph*)}]
        \item 0 \item 1 \item 2 \item 3
    \end{enumerate}

    \item Banyaknya pasangan empat bilangan asli $(a,b,c,d)$ dengan $a<b<c<d$ sehingga himpunan $\{a,b,c,d\}$ memiliki median dan mean $24$ adalah . . . .
    \begin{enumerate}[label={(\Alph*)}]
        \item 253 \item 276 \item 506 \item 552
    \end{enumerate}
\end{enumerate}

\end{document}